\chapter{Constrained Jacobi fields and spectral tests}\label{sec:cmc-jacobi}
A scalar constraint changes the equation defining the nullspace as well as the allowed test functions. We work on a smooth closed connected two-sided CMC boundary $\Sigma=\partial\Omega$ in a Riemannian manifold and put
\[
 q=|A|^2+\operatorname{Ric}(\nu,\nu),\qquad J=\Delta_\Sigma+q,\qquad
 Q(f,g)=\int_\Sigma(\langle\nabla f,\nabla g\rangle-qfg)\,dA.
\]
The source is \cite{BdCE88}, Section 2, pp.~125--129. The scalar mean curvature in that paper is the negative of our $H/n$; its normalized Ricci tensor is multiplied by $n$ in the potential, giving the usual Ricci tensor used here.

\begin{definition}[The constrained Jacobi nullspace]\label{def:constrained-jacobi-nullspace}
Let $\mathcal Z=\{f\in C^\infty(\Sigma):\int_\Sigma f\,dA=0\}$. A constrained Jacobi speed is an $f\in\mathcal Z$ for which $Q(f,g)=0$ for every $g\in\mathcal Z$. An ordinary Jacobi speed instead satisfies $Jf=0$, with no mean condition built into that equation. This is the distinction in \cite{BdCE88}, Definition 2.8, Proposition 2.9 and Remark 2.11. Being a null vector of a bilinear form is stronger than merely satisfying $Q(f,f)=0$ when the form is indefinite.
\end{definition}

\begin{proposition}[The constrained Jacobi equation has a constant right-hand side]\label{prop:constrained-jacobi-constant}
For $f\in\mathcal Z$, the constrained Jacobi condition is equivalent to
\[
 Jf=b\quad\hbox{on }\Sigma
\]
for some real constant $b$. Moreover $b=\operatorname{Area}(\Sigma)^{-1}\int_\Sigma qf\,dA$. This is \cite{BdCE88}, Proposition 2.9.
\end{proposition}
\begin{proof}
Integration by parts gives $Q(f,g)=-\int gJf$. If $Jf=b$, the latter expression vanishes for all $g$ of zero mean. Conversely let $b$ be the average of the smooth function $Jf$ and take $g=Jf-b\in\mathcal Z$. The nullspace condition gives
$0=-\int(Jf-b)Jf=-\int(Jf-b)^2$, hence $Jf=b$ everywhere. Integrating the equation and using $\int\Delta f=0$ gives the formula for $b$. If $q$ is constant, then $b=0$ since $f$ has zero mean; for variable $q$ that deduction is unavailable.
\end{proof}

\begin{lemma}[Zero quadratic form and stability]\label{lem:cmc-zero-form-nullspace}
Suppose $Q\geq0$ on $\mathcal Z$. For $f\in\mathcal Z$, the condition $Q(f,f)=0$ is then equivalent to the constrained Jacobi condition. Without stability, the implication from zero quadratic form need not hold.
\end{lemma}
\begin{proof}
For every $g\in\mathcal Z$ and $t\in\mathbb R$,
\[
 0\leq Q(f+tg,f+tg)=2tQ(f,g)+t^2Q(g,g)
\]
when $Q(f,f)=0$. A nonzero linear coefficient would make the expression negative for a sufficiently small $t$ of one sign. Thus $Q(f,g)=0$ for every $g$. The converse follows by taking $g=f$. For an indefinite form, $B((a,b),(a,b))=a^2-b^2$ has $B((1,1),(1,1))=0$ but $B((1,1),(1,0))=1$. This algebraic example explains why the stability assumption cannot be removed from the inference.
\end{proof}

\begin{proposition}[Tangents to fixed-volume CMC families]\label{prop:cmc-family-jacobi}
Let $\Sigma_t$ be a smooth family of closed CMC boundaries with constant enclosed volume, and write $H_{\Sigma_t}=c(t)$. Its initial normal speed $f$ satisfies
\[
 \int_{\Sigma_0}f\,dA=0,\qquad Jf=-c'(0).
\]
It is therefore a constrained Jacobi speed. This identifies the geometric meaning of the constant in \cite{BdCE88}, equation (2.10).
\end{proposition}
\begin{proof}
Differentiate volume and use the transport formula to obtain zero mean. Proposition \ref{prop:riemannian-mean-curvature-variation} gives the derivative of mean curvature for a normal variation as $-Jf$. A tangential component of velocity contributes the tangential derivative of $H_{\Sigma_0}$, which vanishes because it is constant. Thus $c'(0)=-Jf$ at every point, and Proposition \ref{prop:constrained-jacobi-constant} applies. This is a necessary condition for an actual CMC family; it does not assert that every constrained Jacobi field integrates to such a family.
\end{proof}

\begin{proposition}[Killing fields give zero-mean Jacobi speeds on boundaries]\label{prop:cmc-killing-jacobi}
Let $K$ be a smooth ambient Killing field defined on a neighborhood of the compact region $\Omega$. Then $f=\langle K,\nu\rangle$ on its CMC boundary satisfies $Jf=0$ and $\int_\Sigma f=0$. Hence it belongs to the constrained Jacobi nullspace. The function can be identically zero if $K$ is everywhere tangent to $\Sigma$. Compare \cite{BdCE88}, Proposition 2.12 and the following paragraph.
\end{proposition}
\begin{proof}
The local flow of $K$ consists of isometries, so it preserves the constant mean-curvature value of $\Sigma$. The normal-variation formula, with zero tangential derivative of $H$, gives $Jf=0$. Tracing the Killing equation gives $\operatorname{div}K=0$ in the ambient manifold. The divergence theorem on $\Omega$ therefore gives $\int_\Sigma f=0$. The boundary hypothesis matters for this last step; an arbitrary two-sided hypersurface need not bound an available region on which $K$ is defined.
\end{proof}

\begin{proposition}[The constant-potential stability criterion]\label{prop:cmc-constant-potential-test}
Suppose $q$ is constant on the closed connected hypersurface and let $\lambda_1>0$ be the first positive eigenvalue of $-\Delta_\Sigma$. Then
\[
 \Sigma\text{ is volume-constrained stable}\quad\Longleftrightarrow\quad q\leq\lambda_1.
\]
If $q>0$ and an ambient Killing field has a nonzero normal speed on $\Sigma$, with the hypotheses of Proposition \ref{prop:cmc-killing-jacobi}, this is equivalent to $q=\lambda_1$. This is the precise nontrivial-mode setting of \cite{BdCE88}, Proposition 2.13 and its preceding paragraph.
\end{proposition}
\begin{proof}
The Rayleigh characterization of the first eigenvalue gives $\int|\nabla f|^2\geq\lambda_1\int f^2$ for every zero-mean smooth $f$. Thus $q\leq\lambda_1$ implies nonnegativity of $Q$. If $q>\lambda_1$, a first eigenfunction gives $Q(f,f)=(\lambda_1-q)\int f^2<0$. Existence and smoothness of this eigenfunction use the compact Laplace spectral theorem. For a nonzero Killing normal speed, $Jf=0$ gives $-\Delta f=qf$, and $f$ has zero mean. Hence $q$ is a positive eigenvalue and $q\geq\lambda_1$. Combining the inequalities gives the last equivalence. Without such a mode, equality is not part of the general criterion.
\end{proof}

\begin{proposition}[Index and nullity for a constant potential]\label{prop:cmc-constant-potential-index}
Under the hypotheses of Proposition \ref{prop:cmc-constant-potential-test}, let the eigenvalues of $-\Delta$ be $0=\lambda_0<\lambda_1\leq\lambda_2\leq\cdots$, counted with multiplicity. The constrained index is the number of $\lambda_j<q$ with $j\geq1$. Its constrained nullspace is the span of the eigenfunctions with $\lambda_j=q$, $j\geq1$. If $q>0$, the unrestricted index is the constrained index plus one.
\end{proposition}
\begin{proof}
Use the external compact spectral theorem to expand $f=\sum_j a_j\phi_j$ in an orthonormal Laplace eigenbasis, with $\phi_0$ constant. The second-variation form is diagonal:
\[
 Q(f,f)=\sum_j(\lambda_j-q)a_j^2.
\]
For smooth $f$ this follows from integration by parts and the spectral expansion of the Dirichlet energy. Zero mean removes exactly $a_0$. The negative eigenspace gives a negative subspace of the asserted dimension; any larger subspace intersects its orthogonal complement and hence contains a vector with nonnegative quadratic form. Proposition \ref{prop:constrained-jacobi-constant} gives $Jf=0$ in the constant-potential case, so the constrained nullspace consists exactly of the eigenmodes with eigenvalue $q$. Finally the constant mode is negative precisely when $q>0$. The discreteness and divergence of the Laplace spectrum ensure finiteness of these counts.
\end{proof}
